\subsection{Baseline Mechanism and Its Limitation}
At the high command, B retains 99.95\% of U's mean speed and reduces the largest sector rate from 1516.8 to 536.4 degree/s (Table~\ref{tab:main}). U exceeds the 540 degree/s budget in every high-command trial; B has no violating frames. G rejects an IK target in all eight high-command trials, indicating that an excursion bound without next-window anticipation can leave an unreachable articulated residual. S and D reject targets in all main trials. These are failures of the specified restrictions and workspace, not universal results for those controller classes.

P attains comparable speed with a smaller largest distal-target rate than B. B's 1405.5 degree/s distal peak, although lower than U's 3299.5, still exceeds the sector budget. B uses 94.2 J of absolute joint work and 0.509 m integrated slip, compared with N's 58.1 J and 0.290 m; higher speed alone is not an efficiency result. Overall, 193/268 trials complete and all 75 terminations are rejected IK targets. B completes 68/68 with no sector violations, including the command and parameter cases. The original 1/2/4 ms check changes paired B/N speed by at most 0.00083 m/s; the geometry change gives at most 0.0123 m/s. These local results do not establish general robustness.

\begin{table*}[t]\centering\small
\caption{All-joint limit study at a 0.45 m/s command: four phases, 20 s per trial}
\label{tab:joint}\begin{tabular}{ccccccc}\toprule
 & & Speed (m/s) & $e_v$ (m/s) & Joint RMS ($^\circ$) & Goal lag ($^\circ$) & Clock fraction\\
Clock & Code & Mean [min, max] & Mean & Worst observed & Maximum & Mean\\\midrule
Unlimited & U & 0.287 [0.283, 0.291] & 0.188 & 19.38 & 0.0 & 1.000\\
Unlimited & B & 0.289 [0.285, 0.291] & 0.186 & 19.59 & 0.0 & 1.000\\
Unlimited & P & 0.289 [0.286, 0.291] & 0.186 & 19.45 & 0.0 & 1.000\\
\midrule
H, $\rho=1.0$ & U & 0.220 [0.218, 0.222] & 0.239 & 7.44 & 45.7 & 0.602\\
H, $\rho=1.0$ & B & 0.223 [0.219, 0.226] & 0.236 & 7.36 & 28.7 & 0.604\\
H, $\rho=1.0$ & P & 0.224 [0.220, 0.228] & 0.234 & 7.37 & 19.5 & 0.607\\
\midrule
H, $\rho=0.8$ & U & 0.231 [0.218, 0.244] & 0.232 & 5.58 & 47.5 & 0.566\\
H, $\rho=0.8$ & B & 0.227 [0.211, 0.241] & 0.235 & 5.54 & 30.0 & 0.552\\
H, $\rho=0.8$ & P & 0.228 [0.213, 0.242] & 0.234 & 5.56 & 20.8 & 0.555\\
\midrule
F, $\rho=0.8$ & U & 0.320 [0.315, 0.323] & 0.169 & 6.62 & 64.8 & 1.000\\
F, $\rho=0.8$ & B & 0.316 [0.309, 0.321] & 0.173 & 6.96 & 49.4 & 1.000\\
F, $\rho=0.8$ & P & 0.317 [0.309, 0.324] & 0.171 & 6.96 & 33.8 & 1.000\\
\bottomrule
\end{tabular}
\par\vspace{5pt}\noindent\parbox{.96\textwidth}{\footnotesize All entries complete 4/4 trials. H holds the gait clock until each goal is issued; F keeps the original gait clock. Joint RMS is the maximum, over trials and joints, of issued-target tracking RMS. Goal lag is requested minus issued target. Clock fraction is gait time divided by wall time. $\rho$ scales the joint-group no-load speeds; it is not a hardware confidence bound.}
\end{table*}

\begin{figure*}[t]\centering\includegraphics[width=.97\textwidth]{joint_limits.pdf}
\caption{High-command joint-limit comparisons. Speed markers and whiskers are four-phase means and ranges; tracking and goal-lag panels show worst observed values. H retimes the gait to issue every target; F keeps the original clock. H 0.5 belongs to the separately prescribed stricter-rate suite. All values include startup. The simple periodic P remains competitive.}
\label{fig:joint}\end{figure*}

\subsection{All-Joint Limits and Clock Choice}
All 176 additional trials complete their prescribed horizons. Every one of the 128 bounded trials satisfies both the issued-target and inner-setpoint rate bounds to $10^{-8}$ degree/s numerical tolerance. This verifies the command construction across the tested trajectories. It does not constrain the actual joint response: across the bounded trials, the largest measured joint-rate-to-declared-limit ratio is 1.232.

At the high command, B's mean speed changes from 0.289 m/s without the governor to 0.227 m/s with H at $\rho=0.8$ (Table~\ref{tab:joint}). The worst observed joint tracking RMS decreases from $19.59^\circ$ to $5.54^\circ$, while its gait clock advances at 55.2\% of wall time. With F at the same limits, B achieves 0.316 m/s and $6.96^\circ$ worst joint RMS, but the requested-to-issued goal lag reaches $49.4^\circ$. Thus rate limiting can alter both tracking and propulsion; it does not merely scale speed. F's higher speed is accompanied by a changed target path, not faithful execution of the original trajectory.



P's high-command mean speed is slightly higher than B's at every bounded setting, including 0.228 m/s for H at $\rho=0.8$ and 0.317 m/s for F. Its maximum goal lag is smaller. At $\rho=0.5$, H yields U/B/P speeds of 0.136/0.145/0.146 m/s. At the low command with H at $\rho=0.8$, the corresponding means are 0.106/0.106/0.107 m/s; F gives approximately 0.155 m/s for all three. These observations provide no consistent transport advantage for B over the simple alternatives.

\subsection{Transitions, Longer Runs, and Sensitivity}
Both 24 s command sequences complete for every code and phase. On the forward/stop/reverse sequence, B's mean translational RMSE increases from 0.110 m/s without the governor to 0.133 m/s with H at $\rho=0.8$; P gives 0.111 and 0.133 m/s. On the yaw/curve sequence, bounded B has translational and yaw-rate RMSE of 0.088 m/s and 0.168 rad/s, versus P's 0.084 m/s and 0.168 rad/s. Completion therefore does not imply accurate command tracking.

All twelve 60 s trials complete. At the high command, H at $\rho=0.8$ yields mean speeds of 0.226/0.216/0.216 m/s for U/B/P. Across bounded trials, loaded support falls to one leg and up to 10.55\% of physics samples have fewer than three loaded legs. The scheduled stance count is not a load guarantee.

The constrained B/P timestep check changes paired net speed by up to 0.0111 m/s relative to 2 ms, larger than in the baseline study. This sensitivity prevents transferring the earlier numerical-stability claim to the governor unchanged. Across all 176 new trials, the root-origin velocity integral agrees with displacement speed within 0.000052 m/s. Independent B/P visual replays (Fig.~\ref{fig:replay}) reproduce their original records and are excluded from trial counts.
